\section{Electromagnetic Resilience Threat Model}\label{sec:emr-threat}

The threat model of this report is stated in electromagnetic terms.
Persistent storage is exposed to perturbations whose proximate cause
is electromagnetic energy deposition in the digital substrate. The
model partitions the cause space into two families and defines a
transverse saturation boundary that constrains the hypothesis of the
recovery soundness theorem.

\subsection{Family A: stochastic electromagnetic perturbations}

Family A encompasses perturbations whose statistics are governed by
non-adversarial physical processes: galactic and solar cosmic-ray
ions, terrestrial neutron and alpha background, and thermal noise
contributing to bit-error rates near design margins. The canonical
treatment for digital microelectronics
is~\cite{dodd2003,baumann2005}, with cross-section measurement and
upset rate calculation conventions
in~\cite{petersen1996,heinrich1977}, and the avionics-environment
specification in~\cite{normand1996}. JEDEC~JESD89A~\cite{jedec89a}
codifies measurement and reporting requirements for terrestrial
soft errors. NAND flash storage substrates are governed by their
own family-A statistics including read disturb, retention, and
program-erase cycling effects~\cite{cai2017nand,mielke2017}.

The Family A perturbation operator acts pointwise on each digital
device exposed to the perturbation field. The relevant scalar
parameter is the upset cross-section \(\sigma\), measured in cm\(^2\)
per device, integrated over a fluence \(\Phi\), measured in
particles per cm\(^2\). The expected number of upsets per device
over an exposure interval is the product \(\sigma \cdot \Phi\). At
device, board, and system level, the per-device statistics
aggregate without correlation; this independence assumption
underpins the SEU rate-calculation method of~\cite{heinrich1977}.

For a filesystem state \(s\) and a Family A event \(e\), the
relevant action is bit-level: each upset toggles one bit at a
position determined by the energy deposition site, with
multiplicity governed by the device's MBU pattern. For modern
high-density CMOS substrates, MBU widths of two to eight bits
within a single physical word are reported routinely under
representative test
beams~\cite{baumann2005,jedec89a,cai2017nand,mielke2017}.

\subsection{Family B: adversarial electromagnetic perturbations}

Family B encompasses perturbations whose source is intentional and
whose waveform parameters are chosen to maximize coupling to the
target. Three regimes are treated, in order of increasing peak
intensity: intentional electromagnetic interference (IEMI),
high-power microwave (HPM), and electromagnetic pulse (EMP). The
foundational treatment is~\cite{radasky2004iemi}; the
classification methodology of facilities is
in~\cite{mansson2008iemi}; the standardization framework is
IEC\,61000-2-13 (environment)~\cite{iec61000-2-13} and
61000-4-36 (test methods)~\cite{iec61000-4-36}.

Family B differs from Family A in three respects relevant to
filesystem design: the perturbation field is correlated across
devices on the same board (because the irradiating waveform is
coherent over the target's spatial extent), the perturbation
timing is choosable by the adversary (so the worst case is a
perturbation timed to coincide with a critical kernel
operation), and the per-event symbol-error count can be driven
arbitrarily high within the propagation envelope of the
adversarial waveform.

For a filesystem state \(s\) and a Family B event \(e\), the
filesystem cannot rely on per-region independence of the upset
distribution as it can in Family A. The recovery operator must
therefore treat each region's correction capacity as the
fundamental design unit, and the threat model must explicitly
acknowledge that an adversarial event with sufficient field
strength saturates the correction capacity of any finite-width
code.

The deployment context for Family B is principally civilian and
defense critical infrastructure: this report's design target is
\emph{Op\'erateurs d'Importance Vitale} (OIV) under the French
critical-infrastructure framework, and the EU NIS2 directive's
essential and important entities~\cite{nis2}. The standardization
context for these deployments includes
MIL-STD-882E~\cite{milstd882e} for U.S. defense system safety,
NASA-HDBK-4002A~\cite{nasa-hdbk-4002a} for in-space charging
mitigation, and the ESA ECSS-Q-ST-60-15C~\cite{esa-eccs60} for
European space radiation hardness assurance.

\subsection{Saturation boundary}

The saturation boundary is the transverse threshold at which the
per-subblock symbol-error count of any single Reed--Solomon
subblock exceeds the code's correction capacity. For
\(\mathrm{RS}(n, k)\) with \(n - k\) parity symbols, the
correction capacity is \(t = \lfloor (n-k)/2 \rfloor\). The BEAMFS v2
\textsc{universal\_inline} scheme uses \(\mathrm{RS}(255, 239)\) with \(n - k = 16\)
parity symbols, yielding \(t = 8\) symbol errors per
subblock~\cite{kernelrslib}.

The saturation boundary is stated as a hypothesis on the
perturbation family, not as a property of the operator. A
filesystem cannot, by any algorithm, recover a perturbation that
saturates its code; the appropriate response is fail-closed with
auditable signalling. Theorem v2.1 (\cref{sec:soundness}) covers
the regime within the boundary; Theorem v2.2 covers the regime at
or beyond the boundary, and formally requires the recovery
operator to emit a journal entry flagged with
\texttt{BEAMFS\_RS\_EVENT\_FLAG\_UNCORRECTABLE} at the affected
\((\text{block}, \text{subblock})\) coordinates before returning
fail-closed.

\subsection{Out of scope}

The following classes of perturbation are explicitly outside the
present model. Single-event functional interrupts (SEFI) on the
storage controller firmware are not treated; recovery requires a
controller reset cycle which is below the filesystem layer.
Single-event latch-up (SEL) on the host system is not treated; it
is a system-level fault requiring power cycling. Cumulative total
ionizing dose (TID) and displacement damage (DDD) effects are not
filesystem perturbations in the SEE sense and are governed by
device qualification, not by recovery operators. Side-channel
electromagnetic emanation (TEMPEST and similar) is a
confidentiality concern, not an integrity concern, and is outside
the recovery calculus' scope. Software supply-chain compromise of
the recovery operator's binary is a separate threat model addressed
by signed release artefacts and reproducible builds, not by the
recovery contract.

\subsection{Companion falsification instrument}

The empirical instrument used to falsify v1 Theorem IV.1 and to
test v2 candidates is RadFI~\cite{desbrieres2026radfi}, an
algebraic fault injection operator implemented as a Linux kernel
module. RadFI exposes a parametrized perturbation operator
\(\varepsilon\) acting on the same state space \(\mathcal{S}\) as
the BEAMFS recovery operator \(\mathcal{G}\), with stated
faithfulness properties (RadFI Theorem IV.1, see~\cite{desbrieres2026radfi}).
For the purposes of the present report, RadFI is used in two
modalities. First, it produced the counter-example that retracted
v1 Theorem IV.1, motivating the present revision. Second, it is
the methodology under which any future revision of the present
theorems must itself be exposed to falsification. The present
report's empirical corroboration of v2 (\cref{sec:empirical}) uses
\texttt{dd}-based byte-level disk corruption as a direct
substitute for RadFI's bio-layer perturbation, equivalent under
the SEU model and reproducible without requiring the RadFI
instrument; the kernel-level RadFI hook diagnostic is deferred to
work post-Zenodo.
